% !TeX program = pdflatex
%==============================================================================
%  Lista Teórica 09 --- Métricas para Classificação
%  O gabarito sai de "Lista teorica 08 - gabarito.tex".
%==============================================================================
\providecommand{\gabaritoopt}{}
\documentclass[11pt,a4paper]{article}
\usepackage[\gabaritoopt]{../../recursos/latex/estilo-lista}

\begin{document}
\cabecalholista{08}{Métricas para Classificação}

%------------------------------------------------------------------------------
\begin{exercicio}[uma matriz de confusão, cinco métricas]
Um sistema de detecção de fraude foi aplicado a $10\,000$ transações, das quais
$500$ eram fraudulentas. O resultado:
\begin{center}\small
\renewcommand{\arraystretch}{1.3}
\begin{tabular}{@{}ll cc@{}}
\toprule
 & & \multicolumn{2}{c}{\textbf{previsto}} \\
 & & fraude & normal \\
\midrule
\multirow{2}{*}{\textbf{real}} & fraude & $200$ & $300$ \\
                               & normal & $100$ & $9\,400$ \\
\bottomrule
\end{tabular}
\end{center}
\begin{enumerate}[label=(\alph*),leftmargin=1.1cm]
  \item Calcule acurácia, precisão, \emph{recall} (sensibilidade),
        especificidade e $F_1$.
  \item Calcule a acurácia do classificador que responde \emph{sempre}
        ``normal''. Compare com a do item (a).
  \item Um relatório afirma: ``o sistema tem 96\% de acurácia''. Reescreva essa
        frase de modo que ela informe o leitor em vez de enganá-lo.
  \item Se o banco decidir baixar o corte para capturar mais fraudes, quais das
        cinco métricas do item (a) necessariamente \emph{sobem}, quais
        necessariamente \emph{descem}, e quais podem ir para qualquer lado?
\end{enumerate}
\end{exercicio}

\begin{solucao}
\textbf{(a)} Com $VP=200$, $FN=300$, $FP=100$, $VN=9\,400$:
\begin{center}\small
\renewcommand{\arraystretch}{1.4}
\begin{tabular}{@{}lll@{}}
\toprule
métrica & fórmula & valor \\
\midrule
acurácia        & $(VP+VN)/n = 9\,600/10\,000$      & $0{,}9600$ \\
precisão        & $VP/(VP+FP) = 200/300$            & $0{,}6667$ \\
\emph{recall}   & $VP/(VP+FN) = 200/500$            & $0{,}4000$ \\
especificidade  & $VN/(VN+FP) = 9\,400/9\,500$      & $0{,}9895$ \\
$F_1$           & $2VP/(2VP+FP+FN) = 400/800$       & $0{,}5000$ \\
\bottomrule
\end{tabular}
\end{center}

\smallskip
\textbf{(b)} Ele acerta todas as $9\,500$ transações normais e erra todas as
fraudes: acurácia $=9\,500/10\,000=\mathbf{0{,}9500}$.

O sistema de verdade tem $0{,}9600$. \textbf{Um ponto percentual} separa um
detector que encontra 40\% das fraudes de um programa que não faz nada.

\smallskip
\textbf{(c)} Uma versão honesta seria: \emph{``o sistema detecta 40\% das fraudes
(\emph{recall}); dois em cada três alertas que ele emite são fraudes de verdade
(precisão); e ele marca indevidamente 1\% das transações normais''}.

O que torna a frase original enganosa não é a acurácia estar errada --- ela está
certa --- e sim ela ser \textbf{quase toda determinada pela prevalência}: dos
$9\,600$ acertos, $9\,400$ são transações normais classificadas como normais, o
que qualquer classificador trivial também faria. A informação está nos $600$
casos restantes, e a acurácia os dilui em dez mil.

\smallskip
\textbf{(d)} Baixando o corte, mais transações são chamadas de fraude, então
$VP$ e $FP$ só podem subir e $FN$ e $VN$ só podem descer.
\begin{itemize}[leftmargin=1.4cm]
  \item \textbf{Sobe necessariamente:} o \emph{recall} $VP/(VP+FN)$ --- o
        numerador não desce e o denominador $VP+FN$ é fixo (o total de fraudes).
  \item \textbf{Desce necessariamente:} a especificidade $VN/(VN+FP)$, pelo mesmo
        argumento com o total de normais.
  \item \textbf{Pode ir para qualquer lado:} a precisão (tipicamente desce, mas
        pode subir se os casos recém-incluídos forem majoritariamente fraudes),
        a acurácia (troca $FN$ por $FP$, e o saldo depende de quantos de cada) e
        o $F_1$ (é média harmônica de duas métricas que se movem em direções
        opostas).
\end{itemize}
A moral: \emph{recall} e especificidade sozinhos são triviais de manipular ---
corte $0$ dá \emph{recall} $1$, corte $1$ dá especificidade $1$. Só faz sentido
olhá-los \textbf{aos pares}, e é isso que a curva ROC faz.
\end{solucao}

%------------------------------------------------------------------------------
\begin{exercicio}[o corte que minimiza o custo]
Nem todo erro custa o mesmo. Seja $c_{FP}$ o custo de um falso positivo e
$c_{FN}$ o de um falso negativo, e suponha que o classificador forneça
$p(\x) = \Prob(Y=1\mid \X=\x)$.
\begin{enumerate}[label=(\alph*),leftmargin=1.1cm]
  \item Escreva o custo esperado de decidir $\widehat y=1$ e o de decidir
        $\widehat y=0$, num ponto $\x$ com probabilidade $p$.
  \item Conclua que a decisão ótima é $\widehat y=1$ quando
        $p > \dfrac{c_{FP}}{c_{FP}+c_{FN}}$.
  \item No sistema do Exercício~1, investigar um alerta custa R\$~10 e uma fraude
        não detectada custa R\$~500. Qual o corte ótimo?
  \item O corte padrão do \texttt{scikit-learn} é $0{,}5$. Que razão de custos
        ele supõe implicitamente? Por que isso é quase sempre errado num problema
        desbalanceado?
\end{enumerate}
\end{exercicio}

\begin{solucao}
\textbf{(a)} Se decidimos $\widehat y=1$, erramos quando $Y=0$, o que acontece
com probabilidade $1-p$; o custo esperado é $(1-p)\,c_{FP}$. Se decidimos
$\widehat y=0$, erramos quando $Y=1$, com probabilidade $p$; o custo esperado é
$p\,c_{FN}$. (Acertar custa zero.)

\smallskip
\textbf{(b)} Decidimos $\widehat y=1$ quando o custo esperado dessa decisão é
menor:
\[
  (1-p)\,c_{FP} < p\,c_{FN}
  \iff c_{FP} < p\,(c_{FP}+c_{FN})
  \iff p > \frac{c_{FP}}{c_{FP}+c_{FN}} .
\]
Note que o corte depende \textbf{apenas da razão} entre os custos: multiplicar os
dois por uma constante não muda nada, o que é tranquilizador --- a decisão não
depende da moeda.

\smallskip
\textbf{(c)} Com $c_{FP}=10$ e $c_{FN}=500$:
\[
  p^\ast = \frac{10}{10+500} = \frac{10}{510} \approx \mathbf{0{,}0196}.
\]
Ou seja: vale a pena investigar qualquer transação com \textbf{2\% de chance} de
ser fraude. Faz sentido --- investigar é 50 vezes mais barato que deixar passar.

\smallskip
\textbf{(d)} O corte $0{,}5$ corresponde a
$\tfrac{c_{FP}}{c_{FP}+c_{FN}}=\tfrac12$, isto é $\mathbf{c_{FP}=c_{FN}}$: ele
supõe que os dois erros custam o mesmo.

Isso é quase sempre errado em problema desbalanceado, e por dois motivos que se
somam. Primeiro, o motivo substantivo: a classe rara costuma ser justamente a que
importa (fraude, doença, falha de equipamento), e deixá-la passar custa muito mais
caro que um alarme falso. Segundo, o motivo estatístico: com prevalência de 5\%,
pouquíssimas observações têm $p(\x)>0{,}5$ --- o modelo pode estar ordenando as
transações perfeitamente e ainda assim classificar \textbf{todas} como normais, porque
nenhuma passa do corte.

É por isso que a conduta correta é separar as duas decisões: \textbf{primeiro estime
$p(\x)$ bem, depois escolha o corte}, e escolha-o pelo custo do problema, não
pelo padrão da biblioteca.
\end{solucao}

%------------------------------------------------------------------------------
\begin{exercicio}[AUC na mão]
A AUC tem uma interpretação probabilística exata:
\[
  \mathrm{AUC} = \Prob\big(\text{escore de um positivo} > \text{escore de um negativo}\big),
\]
com sorteio independente e uniforme entre positivos e entre negativos. Considere
cinco observações:
\begin{center}\small
\renewcommand{\arraystretch}{1.2}
\begin{tabular}{@{}lccccc@{}}
\toprule
observação & A & B & C & D & E \\
escore     & $0{,}9$ & $0{,}8$ & $0{,}6$ & $0{,}4$ & $0{,}1$ \\
$y$        & $1$ & $0$ & $1$ & $1$ & $0$ \\
\bottomrule
\end{tabular}
\end{center}
\begin{enumerate}[label=(\alph*),leftmargin=1.1cm]
  \item Liste todos os pares (positivo, negativo) e calcule a AUC diretamente
        pela definição.
  \item Se você somar $0{,}3$ a todos os escores, a AUC muda? E se você elevar
        todos ao quadrado? Justifique.
  \item Qual é a AUC de um classificador que atribui escores completamente ao
        acaso? E de um que ordena \emph{perfeitamente ao contrário}?
  \item Nesta amostra, qual é a maior acurácia atingível escolhendo bem o corte?
        Compare com a AUC e comente.
\end{enumerate}
\end{exercicio}

\begin{solucao}
\textbf{(a)} Positivos: A ($0{,}9$), C ($0{,}6$), D ($0{,}4$). Negativos: B
($0{,}8$), E ($0{,}1$). São $3\times 2=6$ pares:
\begin{center}\small
\begin{tabular}{@{}lcl@{}}
\toprule
par & comparação & positivo ganha? \\
\midrule
A,B & $0{,}9>0{,}8$ & sim \\
A,E & $0{,}9>0{,}1$ & sim \\
C,B & $0{,}6>0{,}8$ & \textbf{não} \\
C,E & $0{,}6>0{,}1$ & sim \\
D,B & $0{,}4>0{,}8$ & \textbf{não} \\
D,E & $0{,}4>0{,}1$ & sim \\
\bottomrule
\end{tabular}
\end{center}
\[
  \mathrm{AUC} = \frac{4}{6} = \frac23 \approx \mathbf{0{,}6667}.
\]

\smallskip
\textbf{(b)} \textbf{Não muda em nenhum dos dois casos.} A AUC depende só das
comparações ``$>$'' entre escores, e qualquer transformação
\textbf{estritamente crescente} preserva todas elas. Somar $0{,}3$ é crescente;
elevar ao quadrado também é, no intervalo $[0,1]$ em que os escores estão.

Essa é a virtude e o defeito da AUC ao mesmo tempo: ela mede \emph{ordenação} e
é cega a \emph{calibração}. Um modelo cujos escores estão todos entre $0{,}49$ e
$0{,}51$ pode ter AUC $1{,}0$ --- e nenhuma das probabilidades ser utilizável.
(Cuidado com o caso não estritamente crescente: elevar ao quadrado escores
\emph{negativos} inverteria comparações e mudaria a AUC.)

\smallskip
\textbf{(c)} Ao acaso: $\mathrm{AUC}=\mathbf{0{,}5}$, porque com escores
independentes dos rótulos a probabilidade de o positivo ganhar é $1/2$ por
simetria. Perfeitamente ao contrário: $\mathrm{AUC}=\mathbf{0}$.

E note que AUC $0$ é ótimo disfarçado: basta inverter o sinal do escore para obter
$1$. Um modelo com AUC bem abaixo de $0{,}5$ não é inútil --- é um bug de sinal.

\smallskip
\textbf{(d)} Percorrendo os cortes possíveis:
\begin{center}\small
\begin{tabular}{@{}lccccc@{}}
\toprule
corte acima de & --- & $0{,}1$ & $0{,}4$ & $0{,}6$ & $0{,}8$ \\
prevê 1 para   & todos & A,B,C,D & A,B,C & A,B & A \\
acertos        & $3$ & $4$ & $3$ & $2$ & $3$ \\
acurácia       & $0{,}6$ & $\mathbf{0{,}8}$ & $0{,}6$ & $0{,}4$ & $0{,}6$ \\
\bottomrule
\end{tabular}
\end{center}
A melhor acurácia é $\mathbf{0{,}8}$, com o corte logo acima de $0{,}1$: os três
positivos são capturados e o único erro é B, um negativo chamado de positivo.

O contraste com a AUC de $0{,}667$ é o ponto do exercício: as duas medem coisas
diferentes. A acurácia avalia \textbf{uma} decisão, tomada num corte específico; a
AUC resume o desempenho \textbf{sobre todos os cortes}, e por isso não depende de
nenhum. Com cinco observações e classes quase equilibradas as duas ainda são
comparáveis; com prevalência de 5\%, como no Exercício~1, a acurácia se cola na
prevalência e a AUC continua informativa.
\end{solucao}

%------------------------------------------------------------------------------
\begin{exercicio}[mesma AUC, calibrações diferentes]
Dois modelos produzem, para as mesmas quatro observações
$y=(0,\,0,\,1,\,1)$, os escores
\[
  p = (0{,}40,\ 0{,}45,\ 0{,}55,\ 0{,}60),
  \qquad
  q = (0{,}05,\ 0{,}10,\ 0{,}90,\ 0{,}95).
\]
\begin{enumerate}[label=(\alph*),leftmargin=1.1cm]
  \item Calcule a AUC dos dois.
  \item Calcule o \textbf{escore de Brier}, $\frac1n\sum_i(p_i-y_i)^2$, dos dois.
  \item Interprete: o que cada uma das duas métricas está medindo, e qual delas
        você usaria para decidir se pode confiar nas probabilidades?
  \item A Aula~08 relata a medição: naive Bayes e logística praticamente empatam
        em AUC ($0{,}776$ contra $0{,}774$), e o Brier os separa ($0{,}1523$
        contra $0{,}0653$). Explique o que esse par de números diz sobre o naive
        Bayes, ligando ao Exercício~4 da Lista Teórica~08.
\end{enumerate}
\end{exercicio}

\begin{solucao}
\textbf{(a)} Nos dois casos os dois positivos têm escores maiores que os dois
negativos, então todos os $2\times2=4$ pares são vencidos pelo positivo:
\[
  \mathrm{AUC}(p) = \mathrm{AUC}(q) = \mathbf{1{,}0}.
\]
Aliás, $q$ é uma função estritamente crescente de $p$, então o item (b) do
Exercício~3 já garantia a igualdade.

\smallskip
\textbf{(b)}
\[
  \mathrm{Brier}(p) = \tfrac14\big[0{,}40^2+0{,}45^2+0{,}45^2+0{,}40^2\big]
   = \tfrac14\big[0{,}16+0{,}2025+0{,}2025+0{,}16\big] = \mathbf{0{,}18125},
\]
\[
  \mathrm{Brier}(q) = \tfrac14\big[0{,}05^2+0{,}10^2+0{,}10^2+0{,}05^2\big]
   = \tfrac14\big[0{,}0025+0{,}01+0{,}01+0{,}0025\big] = \mathbf{0{,}00625}.
\]
O Brier de $p$ é \textbf{29 vezes} o de $q$, com AUC idêntica.

\smallskip
\textbf{(c)} A \textbf{AUC} mede \emph{ordenação}: com que confiabilidade o modelo
coloca positivos acima de negativos. É a métrica certa quando a decisão final
será tomada por um corte que você ainda vai escolher --- e é insensível a
qualquer reescala monótona dos escores.

O \textbf{Brier} é um erro quadrático entre a probabilidade prevista e o rótulo
observado, e mede \emph{calibração} além de ordenação: ele pune tanto ordenar mal
quanto acertar a ordem com probabilidades tímidas ou exageradas. O modelo $p$
ordena perfeitamente, mas todas as suas probabilidades ficam entre $0{,}40$ e
$0{,}60$; se alguém tomar $p_i=0{,}55$ ao pé da letra, vai agir como se houvesse
$45\%$ de chance de erro onde não há nenhuma.

Para decidir se \emph{pode confiar nas probabilidades}, use o \textbf{Brier} (ou
uma curva de calibração). A AUC não responde a essa pergunta, por construção.

E a distinção não é acadêmica: sempre que a probabilidade entra numa conta ---
custo esperado (Exercício~2), preço de um seguro, priorização de uma fila de
espera --- é calibração que importa, não ordenação.

\smallskip
\textbf{(d)} O par de números diz que \textbf{o naive Bayes ordena tão bem quanto
a logística e estima probabilidades muito pior}. A AUC praticamente empatada
($0{,}776$ contra $0{,}774$) mostra que ele coloca os positivos acima dos
negativos com a mesma eficácia; o Brier $2{,}3$ vezes maior ($0{,}1523$ contra
$0{,}0653$) mostra que os valores que ele imprime não são probabilidades
confiáveis.

O mecanismo é o do Exercício~4 da Lista Teórica~08: o naive Bayes multiplica as
densidades marginais como se as covariáveis fossem independentes. Quando elas são
correlacionadas --- e são --- ele conta a \emph{mesma} evidência várias vezes, e
o produto de muitos fatores concordantes empurra a posteriori para $0$ ou $1$. O
resultado são probabilidades sistematicamente extremas.

Mas o excesso de confiança é \textbf{monótono}: ele exagera na mesma direção em
que a evidência aponta, então preserva a ordem. Por isso a AUC não sofre. Numa
frase: \emph{ordenar bem não é estimar bem} --- e escolher a métrica é escolher
qual das duas coisas você precisa.
\end{solucao}


\end{document}
